% =====================================================================
%  CHAPTER 0: 绪论 — 参考：0310营养参考_extracted.txt (PPT原文)
% =====================================================================
\chapter{绪论}
\chapterintro{
\keyword{掌握}营养学与食品卫生学的联系与区别；掌握食品、营养、营养素、膳食、食品安全等基本概念的准确定义；掌握六大营养素分类及宏量营养素与微量营养素的划分；掌握必需营养素的定义及条件必需营养素的概念。\keyword{熟悉}近代营养学关注的七大主要问题；熟悉DRIs七项指标体系；熟悉食品安全的两大关键要素及四大属性。\keyword{了解}食药物质（药食同源物质）的概念及纳入要求。
}

% ----- 第一节 营养学与食品卫生学的联系与区别 -----
\section{营养学与食品卫生学的联系与区别}

营养学与食品卫生学的联系在于，从广义上讲，两者有共同的研究对象：\keyword{食物（food）}和\keyword{人体}，即研究食物（饮食）与健康的关系；区别在于，从狭义上讲，两者在具体研究目标、研究目的、研究方法和理论体系等方面各不相同。\imp{营养学是研究食物中的有益成分与健康的关系}，\imp{食品卫生学则是研究食物中有害因素与健康的关系}。

% ----- 第二节 基本概念 -----
\section{基本概念}

\subsection{食品、食物与食药物质}

\begin{defbox}
\textbf{食品：}指各种供人食用或者饮用的成品和原料以及按照传统既是食品又是中药材（药品）的物品，但是不包括以治疗为目的的物品。\imp{一切食物都是食品}。
\end{defbox}

\begin{keybox}
\textbf{食物（Food）的基本概念与功能：}

\textbf{食物：}供人类食用/饮用；含有营养素；不含有害物质；不是治疗用的药物。

食物的基本功能包括：
\begin{enumerate}[leftmargin=*]
    \item \textbf{生理功能：}提供能量和营养素，维持生命活动
    \item \textbf{心理功能：}满足感官享受和心理需求
    \item \textbf{社会功能：}饮食文化、社交媒介
\end{enumerate}
\end{keybox}

\begin{defbox}
\textbf{食药物质（药食同源物质）：}传统作为食品，且列入《中华人民共和国药典》的物质。纳入食药物质目录的物质应当符合下列要求：
\begin{enumerate}[leftmargin=*]
    \item 有传统上作为食品食用的习惯；
    \item 已经列入《中华人民共和国药典》；
    \item 安全性评估未发现食品安全问题；
    \item 符合中药材资源保护、野生动植物保护、生态保护等相关法律法规规定。
\end{enumerate}
\end{defbox}

\subsection{营养与营养素}

\begin{defbox}
\textbf{营养（Nutrition）：}生物从外界摄入食物，在体内经过消化、吸收、代谢以满足其自身生理功能和从事各种活动需要的必要生物学过程。整个过程包括：\imp{消化吸收 $\rightarrow$ 中间代谢 $\rightarrow$ 排泄}。
\end{defbox}

\begin{defbox}
\textbf{膳食（Diet）：}人们日常食用的饮食，由多种食物组成的。
\end{defbox}

\begin{keybox}
\textbf{营养素（Nutrients）六大分类：}

营养素是指维持机体繁殖、生长发育和生存等一切生命活动和过程，需要从外界环境中摄取的物质。来自食物中的营养素种类繁多，人类所需大约40多种，根据其化学性质和生理作用分为六大类：
\begin{enumerate}[leftmargin=*]
    \item \textbf{蛋白质（Protein）}——提供能量，构建机体，调节代谢
    \item \textbf{脂类（Lipids）}——提供能量，构成细胞膜，溶解脂溶性维生素
    \item \textbf{碳水化合物（Carbohydrate）}——最主要的能量来源
    \item \textbf{矿物质（Mineral）}——构成机体组织，调节生理功能
    \item \textbf{维生素（Vitamin）}——调节机体代谢
    \item \textbf{水（Water）}——构成人体组织，参与物质代谢
\end{enumerate}

\imp{根据人体的需要量或体内含量多少，可将营养素分为宏量营养素和微量营养素：}

\begin{itemize}[leftmargin=*]
    \item \textbf{宏量营养素（Macronutrients）：}人体需要量较多的营养素，包括蛋白质、脂类和碳水化合物。这三种营养素可通过体内的氧化过程释放能量，又称为\imp{产能营养素（Energy-yielding Nutrients）}。
    \item \textbf{微量营养素（Micronutrients）：}相对宏量营养素来说，人体需要量较少的营养素，包括矿物质和维生素。根据在体内的含量不同，矿物质又可分为\imp{常量元素（Macroelements，$>0.01\%$体重）}和\imp{微量元素（Microelements，$<0.01\%$体重）}。根据溶解性质不同，维生素又可分为\imp{脂溶性维生素（Lipid-soluble Vitamins）}和\imp{水溶性维生素（Water-soluble Vitamins）}。
\end{itemize}
\end{keybox}

\begin{keybox}
\textbf{必需营养素与条件必需营养素：}

\textbf{必需营养素（Essential Nutrients）：}机体存活、正常生长和功能所必需，但不能由机体自身合成或合成不足，而必须从食物中获得的营养素。与其他膳食成分相比，缺乏必需营养素可造成特异性缺乏病，甚至死亡。

\textbf{条件必需营养素（Conditionally Essential Nutrients）：}特指人体正常状态下不一定需要，但对于体内不能足量合成的人群是必需供给的营养素。补充该类营养素可纠正缺乏导致的异常表现。这一概念用于全肠外营养的患者、生长发育不全、某些病理状态以及遗传缺陷等条件下人体所需的营养素。

\textbf{其他膳食成分：}食物中除必需营养素以外的具有生物活性的物质，如植物化学物（Phytochemicals），对维护健康、预防疾病具有重要意义。
\end{keybox}

% ----- 第三节 近代营养学主要关注问题 -----
\section{近代营养学主要关注问题}

\begin{keybox}
\textbf{近代营养学关注的七大领域：}
\begin{enumerate}[leftmargin=*]
    \item \textbf{膳食营养素供给量和膳食指南：}营养素的生理功能不仅仅局限于预防营养缺乏病的作用。而营养素发挥这些新功能一般都需要比以往的人体需要量或推荐的供给（\keyword{RDA：供给量}，摄入上下限。\keyword{DRI：参考量}）更高的摄入量。为了满足当前消费者预防慢性病和延缓衰老，而增加营养素摄入量的需求，美国学者首先提出：每日参考摄入量（Daily Reference Intake, DRI）的概念，在RDA的基础上，增加了\imp{适宜摄入量（Adequate Intake, AI）}和\imp{摄入量高限（Upper Limit, UL）}。
    \item \textbf{营养与慢性病：}研究膳食营养因素在肥胖、糖尿病、心血管疾病、癌症等慢性病发生发展中的作用及预防策略。
    \item \textbf{营养与基因关系研究：}营养基因组学（Nutrigenomics）探讨营养素对基因表达的调控及基因多态性对营养素需求的影响。
    \item \textbf{食物成分研究：}系统分析食物中已知和未知的营养成分及生物活性成分。
    \item \textbf{营养与抗氧化功能研究：}在正常生理状态下，人体内的\imp{氧自由基（Oxygen Free Radical, OFR）}不断产生又不断被清除，基本处于动态平衡状态。当机体内氧化和抗氧化之间的平衡严重偏向氧化状态时，称为\imp{氧化应激状态（Oxidative Stress）}。此时，过多的活性氧自由基导致细胞损伤。
    \item \textbf{膳食指南（Dietary Guidelines）：}根据营养学原理提出的一组以食物为基础的建议，指导大众合理选择与搭配食物。
    \item \textbf{功能食品（Functional Food）：}除基本营养功能外，具有调节机体功能、预防疾病作用的食品。
\end{enumerate}
\end{keybox}

% ----- 第四节 食品安全 -----
\section{食品安全的概念及与食品卫生的关系}

\begin{defbox}
\textbf{食品安全（Food Safety）：}
\begin{enumerate}[leftmargin=*]
    \item 对食品按其原定用途进行制作、食用时不会使消费者健康受到损害的一种担保（1996年世界卫生组织）
    \item 食物中有毒、有害物质对人体健康影响的公共卫生问题（WHO）
    \item 食品安全，指食品无毒、无害，符合应当有的营养要求，对人体健康不造成任何急性、亚急性或者慢性危害（《中华人民共和国食品安全法》）
\end{enumerate}
\end{defbox}

食品安全问题涉及两个关键要素：\imp{一是有毒有害物质}，\imp{一是对人体健康的不良影响}。两个要素必须同时存在，才构成食品安全问题。

\begin{keybox}
\textbf{Food Safety与Food Security的区别：}
\begin{itemize}[leftmargin=*]
    \item \textbf{Food Safety（食品安全）：}侧重食品质的方面，即食品无毒、无害，不会对消费者健康造成损害。
    \item \textbf{Food Security（食品安全或粮食安全）：}通常是指食品量的安全，即是否有能力得到或者提供足够的食物或者食品。（注：security在餐馆业是指保护消费者和雇员不受暴虐和犯罪危害。）
\end{itemize}
\end{keybox}

\begin{defbox}
\textbf{营养安全（Nutrition Security）：}人类日常饮食中有足够、平衡且含有人体发育必需营养素的供给，以达到完善的食品安全。
\end{defbox}

\begin{tipbox}
\textbf{隐性饥饿（Hidden Hunger）：}机体因营养不平衡或缺乏某些维生素/必需矿物质，同时其他营养成分过度摄入，产生隐藏性营养需求的饥饿症状。全球约20亿人遭受隐性饥饿，中国约3亿人受影响。70\%的慢性病（糖尿病、心血管病、癌症等）与营养元素摄取不均衡有关。
\end{tipbox}

\begin{keybox}
\textbf{食品安全与食品卫生的关系：}
\begin{enumerate}[leftmargin=*]
    \item 概念从属关系：\imp{食品安全是属概念（上位概念）}，\imp{食品卫生是种概念（下位概念）}
    \item 覆盖范围不同：食品安全涵盖种植、养殖、加工、包装、贮藏、运输、销售、消费全环节；食品卫生通常不包含种植、养殖环节
    \item 侧重点不同：食品安全要求结果安全与过程安全完整统一；食品卫生更加侧重过程安全
\end{enumerate}
\end{keybox}

\textbf{食品安全的四大属性（国际共识）：}
\begin{enumerate}[leftmargin=*]
    \item \textbf{综合概念：}包含食品卫生、食品质量、食品营养等，覆盖食品全流程
    \item \textbf{社会概念：}属社会治理概念，不同国家/发展阶段面临的食品安全问题不同
    \item \textbf{政治概念：}食品安全与企业、政府的责任和承诺紧密关联
    \item \textbf{法律概念：}各国从社会系统工程角度推行食品安全综合立法
\end{enumerate}

% ----- 第五节 DRIs概览 -----
\section{膳食营养素参考摄入量（DRIs）概览}

\begin{keybox}
\textbf{中国居民膳食营养素参考摄入量（Dietary Reference Intakes, DRIs）七项指标体系：}

DRIs是一组每日平均膳食营养素摄入量的参考值，用于评价膳食营养素供给量能否满足人体需要及是否存在过量摄入风险。

DRIs包括以下七项指标：
\begin{enumerate}[leftmargin=*]
    \item \textbf{平均需要量（Estimated Average Requirement, EAR）：}满足群体中50\%个体需要量的摄入水平，是制定RNI的基础。
    \item \textbf{推荐摄入量（Recommended Nutrient Intake, RNI）：}满足群体中97\%～98\%个体需要量的摄入水平，是个体每日摄入该营养素的目标值。
    \item \textbf{适宜摄入量（Adequate Intake, AI）：}当个体需要量资料不足、无法计算EAR时，通过观察或实验获得的健康人群某种营养素的摄入量。
    \item \textbf{可耐受最高摄入量（Tolerable Upper Intake Level, UL）：}平均每日可以摄入营养素的最高限量，对几乎所有个体不产生健康危害。
    \item \textbf{宏量营养素可接受范围（Acceptable Macronutrient Distribution Ranges, AMDR）：}蛋白质、脂肪和碳水化合物理想的摄入量范围（以占总能量\%表示）。
    \item \textbf{预防非传染性慢性病的建议摄入量（Proposed Intakes for Preventing Non-communicable Chronic Diseases, PI-NCD）：}以NCD一级预防为目标的必需营养素每日摄入量。
    \item \textbf{特定建议值（Specific Proposed Levels, SPL）：}以降低成年人膳食相关NCD风险为目标的其他膳食成分（如植物化学物）的每日摄入量。
\end{enumerate}

注：DRIs每4～5年修订一次。具体指标体系和应用将在第五章"公共营养"中详述。
\end{keybox}

% ----- 第六节 发展历史 -----
\section{营养与食品卫生学的发展历史}

\subsection{营养学发展简史}
\begin{enumerate}[leftmargin=*]
    \item 古代：经验性营养知识积累（如《黄帝内经》中的食养理论）
    \item 18-19世纪：拉瓦锡提出呼吸是氧化过程，能量代谢研究起步
    \item 20世纪：维生素的陆续发现、必需氨基酸概念的提出；中国营养学会成立（1945年）
    \item 21世纪：分子营养学、营养基因组学等新兴领域发展
\end{enumerate}

\subsection{食品卫生学发展历史}
\begin{enumerate}[leftmargin=*]
    \item 周朝（约3000年前）：设置"凌人"一职，专职负责食品冷藏、防腐
    \item 唐代《唐律》：变质肉食致人患病，剩余肉食必须销毁，违者杖九十
    \item 1965年：颁布《食品卫生管理试行条例》
    \item 1979年：颁布《中华人民共和国食品卫生管理条例》，正式步入法制化阶段
    \item 2015年：修订《食品安全法》正式实施
    \item 2018年：完成《食品安全法》修正，组建国家市场监督管理总局
\end{enumerate}

% ----- 第七节 研究内容与方法 -----
\section{研究内容与方法}

\textbf{研究内容：}
\begin{enumerate}[leftmargin=*]
    \item \textbf{营养学部分：}食物营养素和生物活性成分、各类食物营养、特殊人群营养、公共营养、营养与疾病
    \item \textbf{食品卫生学部分：}食品污染及预防、食品添加剂、各类食品卫生、食源性疾病、食品安全风险分析、食品安全监督管理
\end{enumerate}

\begin{tipbox}
\textbf{公众认知 vs 科学标准的食品危害性排序：}

公众排序：食品添加剂 > 环境污染 > 营养素搭配不良 > 天然毒素 > 微生物污染

科学排序：营养素搭配不良（膳食行为）> 微生物污染 > 天然毒素 > 环境污染 > 食品添加剂
\end{tipbox}

% ----- 复习题 -----
\reviewcontent{
\section{复习题}

\subsection*{一、名词解释}
\begin{enumerate}[leftmargin=*, itemsep=8pt]
    \item \textbf{营养（Nutrition）}
    \item \textbf{营养素（Nutrients）}
    \item \textbf{食品卫生学（Food Hygiene）}
    \item \textbf{食品安全（Food Safety）}
    \item \textbf{粮食安全（Food Security）}
    \item \textbf{隐性饥饿（Hidden Hunger）}
    \item \textbf{必需营养素（Essential Nutrients）}
    \item \textbf{条件必需营养素（Conditionally Essential Nutrients）}
    \item \textbf{膳食（Diet）}
    \item \textbf{食药物质（药食同源物质）}
\end{enumerate}

\subsection*{二、简答题}
\begin{enumerate}[leftmargin=*, itemsep=8pt]
    \item \textbf{简述营养学与食品卫生学的联系与区别。}
    \item \textbf{简述营养素的六大分类，并说明宏量营养素与微量营养素的划分。}
    \item \textbf{简述必需营养素与条件必需营养素的定义及区别。}
    \item \textbf{近代营养学主要关注哪些问题？（答出七个方面）}
    \item \textbf{简述食品安全与食品卫生的关系。}
    \item \textbf{简述食品安全的四大属性。}
    \item \textbf{简述DRIs的七项指标体系。}
\end{enumerate}

\subsection*{参考答案要点}

\begin{enumerate}[leftmargin=*, itemsep=5pt]
    \item 见本章各概念定义框。
    \item 联系：共同研究对象——食物和人体，即研究食物与健康的关系。区别：营养学研究食物中的有益成分与健康的关系，食品卫生学研究食物中有害因素与健康的关系。
    \item 六大类：蛋白质（Protein）、脂类（Lipids）、碳水化合物（Carbohydrate）、矿物质（Mineral）、维生素（Vitamin）、水（Water）。宏量营养素（需量多）：蛋白质、脂类、碳水化合物，三者可通过氧化释放能量，又称产能营养素。微量营养素（需量少）：矿物质（又分常量元素$>0.01\%$体重和微量元素$<0.01\%$体重）和维生素（又分脂溶性和水溶性）。
    \item 必需营养素：机体存活、正常生长和功能所必需，但不能由机体自身合成或合成不足，必须从食物中获得的营养素，缺乏可造成特异性缺乏病甚至死亡。条件必需营养素：正常状态下不一定需要，但体内不能足量合成的人群必需供给，适用于全肠外营养、生长发育不全、某些病理状态及遗传缺陷等条件。
    \item ①膳食营养素供给量和膳食指南（从RDA到DRI的转变，增加了AI和UL）；②营养与慢性病；③营养与基因关系研究；④食物成分研究；⑤营养与抗氧化功能研究（氧化应激）；⑥膳食指南（Dietary Guidelines）；⑦功能食品。
    \item 食品安全是属概念（上位概念），食品卫生是种概念（下位概念）。食品安全涵盖种植、养殖、加工、包装、贮藏、运输、销售、消费全环节，要求结果安全与过程安全完整统一；食品卫生通常不包含种植、养殖环节，侧重过程安全。
    \item 综合概念、社会概念、政治概念、法律概念。
    \item 平均需要量（EAR）、推荐摄入量（RNI）、适宜摄入量（AI）、可耐受最高摄入量（UL）、宏量营养素可接受范围（AMDR）、预防非传染性慢性病的建议摄入量（PI-NCD）、特定建议值（SPL）。
\end{enumerate}

\newpage
}
